
% Test newcommand inside newcommand for local variable.
% Refs:
%   Inner command without arguments:
%   https://tex.stackexchange.com/questions/ [... append next line for full url]
%       28704/defining-a-newcommand-with-variable-name-inside-another-newcommand
%   Inner command with arguments:
%   https://tex.stackexchange.com/questions/ [... append next line for full url]
%       34697/variable-name-newcommand-with-parameters-within-another-newcommand


\documentclass[11pt,letterpaper]{article}

% Option 'draft' does the following:
%   1. Sets the number of TOC levels to 5 from the default 3
%   2. Displays hidden text in \devNotes commands
% Color options:
%       'nocolor' disables non-hyperref color (default)
%    'allowcolor'   allows non-hyperref color
% Options controlling sections starting on new pages:
%        'sepsec'          starts sections on new pages (default)
%      'nosepsec' does not start  sections on new pages
%\usepackage[draft,allowcolor]{latex/gfr}
%\usepackage[draft,allowcolor]{latex/gfr}
%\usepackage[draft]{latex/gfr}
 \usepackage[allowcolor]{latex/gfr}
%\usepackage{latex/gfr}

% Override section number depth.
% \setcounter{tocdepth}{5}
% \setcounter{secnumdepth}{5}

\usepackage{geometry}    % page layout functions
% Page Margins
% Omit textwidth to prevent warning "over-specification in h-direction".
%   textwidth = pagewidth - left - right = 8.50 - 1.00 - 1.00 = 6.50
% \geometry{textwidth     = 6.50in} % computed = pagewidth - left - right
  \geometry{textheight    = 9.00in}
  \geometry{columnsep     = 0.25in}
  \geometry{left          = 1.00in}
  \geometry{right         = 1.00in}
  \geometry{top           = 0.73in} % = 1.00in-0.10in-12pt
  \geometry{bindingoffset = 0.00in, twoside}
  \geometry{headsep       = 0.10in}
  \geometry{footskip      = 0.50in}

\usepackage{changepage}
  % \begin{adjustwidth}{<left_indent>}{<right_indent>}

% Package marginnote was added as a required package in gfr.sty because it is
% used in the debugging versions of FINALFIX which is also defined in gfr.sty.
%\usepackage{marginnote}  % \marginnote{} that does not float.

\usepackage{tocloft}     % to configure the table of content (ToC)
% \setlength{   \cftsecnumwidth}{4.5ex}% width of    section number in ToC
  \setlength{\cftsubsecnumwidth}{6.5ex}% width of subsection number in ToC

\usepackage{multicol}    % table entries spanning multiple columns
\usepackage{multirow}    % table entries spanning multiple row

\usepackage{varwidth}    % minimal width minipage
  % Provides:  \begin{varwidth}{\textwidth}

\usepackage{overcite}    % superscript citations

\usepackage{graphicx}    % image inclusion
  \graphicspath{{Figures}}

\usepackage{subcaption}  % subfigure environment

\usepackage[inline]{enumitem}    % Better control of lists.
% Change description environment to item at top and description indented below.
% Ref:  https://www.latex4technics.com/?note=KFO
% Initiate with command:
% \begin{description}[align=myparleft]
% \end{description}
    \makeatletter
    \SetLabelAlign{myparleft}{\parbox[t]\textwidth{#1\par\mbox{}}}
    \makeatother

\usepackage{array} % formats for tabular columns
% Define some new left/right/center fixed-width column types. Arg #1=width.
  \newcolumntype{L}[1]{>{\raggedright\let\newline\\%
                         \arraybackslash\hspace{0pt}}m{#1}}
  \newcolumntype{C}[1]{>{\centering\let\newline\\%
                         \arraybackslash\hspace{0pt}}m{#1}}
  \newcolumntype{R}[1]{>{\raggedleft\let\newline\\%
                         \arraybackslash\hspace{0pt}}m{#1}}

\usepackage{threeparttable} % threeparttable environment for tables w/ footnotes
\usepackage{longtable}   % page breaking tables
  % \begin{longtable}...

\usepackage{amsmath}     % improved equation typesetting, gather environments
\usepackage{amssymb}     % math symbols
  % \checkmark
\usepackage{bm}          % bold math symbols with \bm{} commmand
\usepackage{yhmath}      % wide math accents (hat, tilde, etc)
\usepackage{mathtools}   % used for \coloneq symbol
\usepackage{xfrac}

% Math script fonts, provides \mathscr{}. Choose only one package:
%   Often used inside \volume command.
% \usepackage{mathrsfs}           % provides \mathscr{}, narrow spindly V
  \usepackage[mathscr]{eucal}     % provides \mathscr{}, softer curved V
% \usepackage[mathcal]{eucal}     % provides \mathcal{}
% \usepackage{eucal}              % provides \mathcal{}
% \usepackage{euscript}           % provides \EuScript{}
% \usepackage[mathcal]{euscript}  % provides \mathcal{}

\usepackage{latex/turbulence}  % turbulence variable shortcuts

%Package siunitx is loaded by gfr.sty as a required package
%\usepackage{siunitx}
%\sisetup{per-mode = symbol,%
%         display-per-mode = fraction,%
%         inline-per-mode = symbol}
  %Example usage (see gfr.sty for more): \unit{N.s/m^2}
  % Provides \unit{} command to typeset units.
  % Provides S tabular column format to left-align numbers and center them.

% xcolor is a superset of color package
%\usepackage[usenames,dvipsnames]{color}       % color text
\usepackage[dvipsnames]{xcolor}                % color
% Available Colors:
%   GreenYellow, Yellow, Goldenrod, Dandelion, Apricot, Peach, Melon, 
%   YellowOrange, Orange, BurntOrange, Bittersweet, RedOrange, Mahogany, 
%   Maroon, BrickRed, Red, OrangeRed, RubineRed, WildStrawberry, Salmon, 
%   CarnationPink, Magenta, VioletRed, Rhodamine, Mulberry, RedViolet, 
%   Fuchsia, Lavender, Thistle, Orchid, DarkOrchid, Purple, Plum, Violet, 
%   RoyalPurple, BlueViolet, Periwinkle, CadetBlue, CornflowerBlue, 
%   MidnightBlue, NavyBlue, RoyalBlue, Blue, Cerulean, Cyan, ProcessBlue, 
%   SkyBlue, Turquoise, TealBlue, Aquamarine, BlueGreen, Emerald, 
%   JungleGreen, SeaGreen, Green, ForestGreen, PineGreen, LimeGreen, 
%   YellowGreen, SpringGreen, OliveGreen, RawSienna, Sepia, Brown, Tan, Gray 
% Available Macros:
%   \colorbox{back_color}{text}
%   \fcolorbox{frame_color}{back_color}{text}
% User Macros:
  \definecolor{nmlcolor}{gray}{0.95}

\usepackage{pdfcolmk}    % fixes color across page breaks

\usepackage{xstring}
  % \IfEq
  % \IfEqCase

\usepackage{calc}        % for simple arithmetic
  % \widthof{}

%\usepackage{tikz}        % drawing macros
%  \usetikzlibrary{math}
%  \usetikzlibrary{calc}  % calculations
%  \usetikzlibrary{arrows}  % arrowheads
%  \usetikzlibrary{arrows.meta}  % arrowheads
%  \usetikzlibrary{decorations.markings}
%% \usetikzlibrary{positioning}
%% \usetikzlibrary{backgrounds}
%\usepackage{tikz-3dplot} % drawing macros 3d
%\usepackage {pgfplots}

\usepackage{listings}    % Source code listings.
% Old style for namelist blocks. Renamed it to clarify it is a lststyle.
% \lstdefinestyle{nml}{%
%    frame=single %
%   ,backgroundcolor=\color{nmlcolor} %
%   ,captionpos=b %
%   ,language={Fortran} %
%   }
  %
  % Style to typeset a namelist block listing.
  \lstdefinestyle{lststylenmlblock}{%
     language=[90]Fortran
    ,commentstyle=\small\ttfamily % same as basicstyle
    ,basicstyle=\ttfamily
    ,columns=fixed
    ,basewidth={0.52em,0.45em} % default: {0.6em,0.45em}
%   ,basewidth={0.50em,0.45em} % default: {0.6em,0.45em}
      % Set the basewidth for a single character, for {fixed,flexible} modes.
      % Use fixed=0.50em to match \texttt, but underscores are different.
    ,frame=single
    ,backgroundcolor=\color{nmlcolor}
    ,captionpos=b
    }
  %
  % Environment to typeset a namelist block listing.
  %   \begin{lstnmlblock}[options]
  %   \end{lstnmlblock}
  \lstnewenvironment{nmlblocklst}[1][]
    {\lstset{style=lststylenmlblock,#1}}
    {}

\lstset{morecomment=[is]{\%@\#!}{!\#@\%}}

\usepackage[most]{tcolorbox}  % Color boxes
  % \tcblisting
  % loads \usepackage{listings}
  %   \lstlisting from \usepackage{listings}
  %   \lstinline  from \usepackage{listings}
  %
  % Environment to typeset a namelist item listing.
  %   \begin{nmldescitem}[options]
  %   \end{nmldescitem}
\tcbuselibrary{listings}
  \newtcblisting{nmldescitem}[1][]{%
     listing only      % Show source listing, but not compiled text.
%-- Autofitted width settings. Default fit is tight, so add some spacings.
    ,hbox              % fit box to content width
    ,left   = 4.0pt    % l padding between content and frame
    ,right  = 4.0pt    % r padding between content and frame
    ,top    = 0.0pt    % t padding between content and frame
    ,bottom = 0.0pt    % b padding between content and frame
    ,boxsep = 1.0pt    % additional lrtb padding between content and frame
%-- Frame Settings
    ,boxrule       = 0.4pt    % thickness of the frame rule on all edges
    ,sharp corners = all      % square corners
    ,colframe      = black    % frame color
    ,colback       = nmlcolor % background color
%-- Spacing Outside Frame
%   ,aboveskip  = 0.0pt % default: \medskipamount
%   ,belowskip  = 0.0pt % default: \medskipamount
    ,after skip = 5.0pt % default: ???
%   ,listing options={belowskip = 0.0pt}  % default: \medskipamount
%-- Font Settings (Style 1)
    ,listing options={language={[90]Fortran}}
    ,fontupper=\ttfamily % \textt font
    ,listing options={columns=fixed}
      % columns=fixed,flexible,spaceflexible,fullflexible
    ,listing options={basewidth={0.52em,0.45em}} % default: {0.6em,0.45em}
      % Set the basewidth for a single character, for {fixed,flexible} modes.
%-- Font Settings (Style 2)
%   ,fontupper=\ttfamily % \textt font
%   ,listing options={flexiblecolumns=true}  % like \texttt in normal text
%-- Escape sequences to access Latex commands
%   ,listing options={escapechar=@}
%   ,listing options={escapeinside={!@}{!@}}
%-- Other Settings
    ,#1                % other optional arguments
    }

\usepackage{nicematrix}  % array/matrix with headers
  \NiceMatrixOptions{code-for-first-row = \color{blue},
                     code-for-first-col = \color{blue},
                     code-for-last-row  = \color{blue},
                     code-for-last-col  = \color{blue}}
  % \begin{bNiceArray}

%\usepackage[acronym,index,toc]{glossaries}  % makeglossaries
%% Could not get acronyms to appear properly in text and index.
%%
%% Reset acronyms to use full form each chapter or section.
%% \preto\chapter{\glsresetall}
%% \preto\section{\glsresetall}
%% Disable \makeglossaries if no glossaries used.
%% \makeglossaries
%  \setacronymstyle{long-short}
%% \glsdisablehyper
%  \newacronym{DG}{DG}{Discontinuous Galerkin}
%  \newacronym{FR}{FR}{Flux Reconstruction}
%  \newacronym{RK}{RK}{Runge-Kutta}
%  \newacronym{CFD}{CFD}{computational fluid dynamics}
%  \newacronym{CFL}{CFL}{Courant-Friedrichs-Lewy}
%  \newacronym{DNS}{DNS}{Direct Numerical Simulation}
%  \newacronym{DoF}{DoF}{degrees of freedom}
%  \newacronym{LES}{LES}{Large-Eddy Simulation}

\usepackage{imakeidx}             % makeindex
  \makeindex[columns=2, title=Index, intoc, options=-r]
    % intoc : add index to table of contents
    % Option -r : supresses adjacent page refs from being grouped together.
  % Use \index{entry|primeidx} to format entry as the primary reference.
    \newcommand{\primeidx}[1]{#1\textsuperscript{*}}
%   \newcommand{\primeidx}[1]{\textbf{#1}}
  % Commands:
  %   \index{name}
  %   \index{name!subname}
  %   \index{name|primeidx}        % mark as a prime entry
  %   \index{name|see{other name}} % alias entry in index

\usepackage[verbose,hyperpageref]{backref}  % back refs from citations
  \renewcommand*{\backreflastsep}{, }  % default={, and }
  \renewcommand*{\backreftwosep}{, }   % default={ and }
  \renewcommand*{\backref}[1]{}
  \renewcommand*{\backrefalt}[4]{%
    % #1 = backref count w/o duplicates
    % #2 = backref list  w/o duplicates
    % #3 = backref count w   duplicates
    % #4 = backref list  w   duplicates
    \ifcase #1 %
%     No citations.% use \relax if you do not want the "No citations." message
      \relax%
    \or
%     This section is for one citation.
      (Cited on page #2).%
    \else
%     This section is for multiple citations.
      (Cited on pages #2).%
    \fi%
  }

% Note:
%   Load imakeidx before hyperref.
%   Load cleveref after  hyperref.
\usepackage[pagebackref]{hyperref}  % hyperref links within pdf
  % Options:
  %   pagebackref  Provides back refs for citations.
  % Commands:
  %   \hypertarget{name}{text}  % Create target, other than label.
  %   \hyperlink{name}{text}    % Point to hyperlink name.
  %   \hyperref[label]{text}    % Point to label.
  %   \nameref{label}           % Point to label, use name as text.
  \hypersetup{           % replace hyperref boxes with colored links
    colorlinks,
    linkcolor={red!50!black},
    citecolor={blue!50!black},
    urlcolor={blue!80!black}
  }

\usepackage{cleveref}    % \Cref{}, \cref{} upper/lowercase clever refs
%% Usage:
%%   \Cref{} is for beg of sentence.
%%   \cref{} is for mid of sentence.
%% \crefname{type}{singular}{plural} Defines the name-prefix for the ref.
%  \Crefname{equation}{Eq.}{Eqs.}
   \crefname{equation}{Eq.}{Eqs.}
   \crefname{figure}{Fig.}{Figs.}
   \Crefname{paragraph}{Section}{Sections}
   \crefname{paragraph}{section}{sections}
   \Crefname{subparagraph}{Section}{Sections}
   \crefname{subparagraph}{section}{sections}
   \Crefname{appendix}{Appendix}{Appendices}
   \crefname{appendix}{appendix}{appendices}
%% \crefrormat{type}{format}
%%   #1 = the formatted label counter
%%   #2 = beg mark for hyperref link
%%   #3 = end mark for hyperref link
   \Crefformat{equation}{Eq.~(#2#1#3)}
   \crefformat{equation}{Eq.~(#2#1#3)}
   \Crefformat{section}{Section~#2#1#3}
   \crefformat{section}{Section~#2#1#3}
   \Crefformat{paragraph}{Section~#2#1#3}
   \crefformat{paragraph}{Section~#2#1#3}
   \Crefformat{subparagraph}{Section~#2#1#3}
   \crefformat{subparagraph}{Section~#2#1#3}
   \crefformat{appendix}{Appendix~#2#1#3}
   \Crefformat{appendix}{Appendix~#2#1#3}
   \newcommand{\crefrangeconjunction}{--} % replace "to" with dash in ranges

\usepackage{comment}     % comment out blocks of text
% Do not use \includecomment{MyComment} or color will be lost.
% Instead, simply disable the \excludecomment{MyComment} line.
  \specialcomment{MyComment}{\begingroup\color{Red}}{\endgroup}
  \excludecomment{MyComment}

\usepackage{booktabs} % \toprule \bottomrule

% algorithm and algorithmic are both part of the algorithms bundle
\usepackage{algorithm}   % provides algorithm environment
\usepackage{algorithmic} % provides algorithmic environment

% Provides \nth{#} command which automates the text exponent for a number
% e.g., 1^st, 2^nd, 3^rd, 4^th, etc.
\usepackage[super]{nth}


% Modify default settings for overcite package.
% References list numbered 1. rather than [1]
  \makeatletter \renewcommand\@biblabel[1]{#1.} \makeatother

% ===========
% Text Macros
% ===========

% Cite in the line of text.
% \newcommand{\Linecite}[1]{Ref.~\citen{#1}} % Start of sentence.
% \newcommand{\linecite}[1]{Ref.~\citen{#1}} % Elsewhere in sentence.
  % Line cite at start of sentence.
  \newcommand{\Linecite}[1]{%
    \def\tempx{0}%
    \StrCount{#1}{,}[\tempx]% Requires \usepackage{xstring}
    Ref%
    \ifnum\tempx > 0 % Add 's' if multiple refs.
      s.%
    \fi%
    ~\citen{#1}%
  }
  % Line cite elsewhere in sentence.
  \newcommand{\linecite}[1]{%
    \def\tempx{0}%
    \StrCount{#1}{,}[\tempx]% Requires \usepackage{xstring}
    Ref%
    \ifnum\tempx > 0 % Add 's' if multiple refs.
      s.%
    \fi%
    ~\citen{#1}%
  }

% Cross-reference commands in place of \ref.
  \newcommand{\eref}[1]{\cref{#1}} % lower equation
  \newcommand{\Eref}[1]{\Cref{#1}} % upper Equation
  \newcommand{\fref}[1]{\cref{#1}} % lower figure
  \newcommand{\Fref}[1]{\Cref{#1}} % upper Figure
  \newcommand{\lref}[1]{\cref{#1}} % lower listing
  \newcommand{\Lref}[1]{\Cref{#1}} % upper Listing
  \newcommand{\sref}[1]{\cref{#1}} % lower section
  \newcommand{\Sref}[1]{\Cref{#1}} % upper Section
  \newcommand{\tref}[1]{\cref{#1}} % lower table
  \newcommand{\Tref}[1]{\Cref{#1}} % upper Table
  \newcommand{\aref}[1]{Appendix~\ref{#1}} % lower appendix
  \newcommand{\Aref}[1]{Appendix~\ref{#1}} % upper Appendix
  \newcommand{\pref}[1]{page~{#1}} % lower page
  \newcommand{\Pref}[1]{Page~{#1}} % upper Page

% For consistency:
  \newcommand{\NOTE}{\textbf{NOTE: }}
  \newcommand{\EXAMPLE}{\textbf{EXAMPLE: }}
  \newcommand{\supst}{\textsuperscript{st}}  % 1^{st}
  \newcommand{\supnd}{\textsuperscript{nd}}  % 2^{nd}
  \newcommand{\suprd}{\textsuperscript{rd}}  % 3^{rd}
  \newcommand{\supth}{\textsuperscript{th}}  % 4^{th}
  %> Format for CGNS output variable names.
  \newcommand{\ov}[1]{\textbf{\textsf{#1}}}

% Command for typesetting "See Also:" with additional info.
% Spacing control requires \usepackage{enumitem}.
\newcommand{\SeeAlso}[1]{%
  \begin{description}[leftmargin=\widthof{See Also:}+1.5em,
                      rightmargin=0.0em,
                      style=sameline]
    \item[See Also:]
    \begin{sloppypar}
    \raggedright
    #1
    \end{sloppypar}
  \end{description}
}

%%%%%%%%%%%%%%%%%%%%%%%%%%%%%%%%%%%%%%%%%%%%%%%%%%%%%%%%%%%%%%%%%%%%%%%%%%%

% \descitem{} macro replacement for \item[] in a \description.
% Allows \lstinline to be used inside \descitem{}.
% Ref:  https://comp.text.tex.narkive.com/dzHkwJ7t/lstinline-in-item
%
\newsavebox{\itembox}
\newcommand*{\descitem}{%
\everyhbox{\everyhbox{}\aftergroup\finishdescitem}%
\setbox\itembox\hbox
}
\newcommand*{\finishdescitem}{%
\everyhbox{}%
\item[\box\itembox]%
}

%##########################################################################
%##########################################################################
%##########################################################################
%##########################################################################
%##########################################################################
%  "Key" commands:
%   These are commands for creating index references that are not for
%   an input variable. References placed in the index from the "key"
%   commands get the default font formatting within the index.
%##########################################################################
%##########################################################################
%##########################################################################
%##########################################################################
%##########################################################################

%#begfold#{{{ Details for command \keyindex
% \keyindex*{#2}
%   #1 = * (or not): prime index entry, else standard index entry
%   #2 = (escaped) text to use for index entry
% Called by other commands to create an input entry within the index.
% Examples:
%   \keyindex*{important phrase}
%      => \index{important phrase|primeidx}
%   \keyindex{important phrase}
%      => \index{important phrase}
% Note:
%   Do not use \keyindex directly, but rather use \key command below.
%#endfold#}}}
%
\NewDocumentCommand\keyindex{sm}{%
  \IfBooleanTF#1
    {\index{#2|primeidx}}% #1=star is present
    {\index{#2}}%          #1=star is missing
}

%%%%%%%%%%%%%%%%%%%%%%%%%%%%%%%%%%%%%%%%%%%%%%%%%%%%%%%%%%%%%%%%%%%%%%%%%%%

%#begfold#{{{ Details for command \key
% \key*[#2]{#3}[#4]    # Define/Reference key text.
%   #1 = * (or not): prime index entry, else standard index entry
%   #2 = hypertarget  (default=none)
%   #3 = (escaped) text to use within paragraph
%   #4 = (escaped) text to use within index entry (default={})
%        If #4 is {=}, then default to #3.
% - Used to define/reference keywords within the document.
% - Typeset the keyword.
% - Create/reference a hyperlink target.
%   - If hyperlink not wanted, skip optional arg #2.
% - Add the keyword to the index.
% Note:
%   Best to avoid special symbols in the hypertarget.
%   hypertarget with \_ gives undefined control sequence \textunderscore.
%   hypertarget with  _ is ok, but still best to avoid.
%#endfold#}}}
%
\NewDocumentCommand\key{sO{}mO{}}{%
  \IfBooleanTF#1{% #1=star is present
    % Typeset the keyword text in the paragraph, create optional hypertarget.
    \IfEq{#2}{}% Test if #2=hypertag is missing.
      {#3}%                   #2 is missing
      {\hypertarget{#2}{#3}%
       \phantomsection\label{#2}%
      }% #2 is present
      % pageref works with labels but not hypertargets.
      % phantomsection creates a local anchor, not the previous section header.
    % Typeset the keyword text in the index as a primary entry.
    \IfEqCase{#4}{%
      {}{}%                 If #4=''  index nothing.
      {=}{\keyindex*{#3}}%  If #4='=' index #3. (ie, copy the paragraph text)
    }[\keyindex*{#4}]%      Else      index #4.
  }{% #1=star is missing
    % Typeset the input variable, create optional hyperlink.
    \IfEq{#2}{}% Test if #2=hypertag is missing.
      {#3}%                 #2 is missing
      {\hyperlink{#2}{#3}}% #2 is present
    % Typeset the keyword text in the index.
    \IfEqCase{#4}{%
      {}{}%                 If #4=''  index nothing.
      {=}{\keyindex{#3}}%   If #4='=' index #3. (ie, copy the paragraph text)
    }[\keyindex{#4}]%       Else      index #4.
  }%
}

%##########################################################################
%##########################################################################
%##########################################################################
%##########################################################################
%##########################################################################
%  "inp" commands:
%   These are commands for creating index references that are for
%   an input variable EXCEPT for any variables within the BC_Input
%   derived type. References placed in the index from the "inp"
%   commands get special formatting/font styling based on the /inpstyle
%   \inpstyle command before placement in the index.
%##########################################################################
%##########################################################################
%##########################################################################
%##########################################################################
%##########################################################################

%#begfold#{{{ Details for command \inpstyle
% \inpstyle[#1]{#2}
%   #1 = (escaped) pretext to typeset before #2
%   #2 = (escaped) text to typeset as an input variable.
% Called by other commands to apply the style for an input variable.
% Do not call directly.  Use \inptext instead.
%#endfold#}}}
%
\NewDocumentCommand\inpstyle{O{}m}{%
%#begfold#{{{ Comment on error when used inside math mode
%% \IfEq gives error:  "incomplete \iffalse" when called from inside math mode
%% with text "var\_name".  Why?
%  \IfEq{#2}{}% Test if #2=text is missing.
%    {% #2=text is missing
%     % do nothing
%    }%
%    {% #2=text is present
%%    \texttt{\MakeLowercase{#1#2}}%
%%    \texttt{\MakeUppercase{#1#2}}%
%     \texttt{{#1}#2}%
%%    {#1#2}%
%%    }%
%#endfold#}}}
%
%% This runs but manditory arg ALWAYS has a value, so \IfValue{#2}=T.
%% Value could be empty, but it has a value.
  \IfValueT{#2}%
    {% #2=text is present
%    \texttt{\MakeLowercase{#1#2}}%
%    \texttt{\MakeUppercase{#1#2}}%
     \texttt{#1#2}%
%    {#1#2}%
    }%
}

%%%%%%%%%%%%%%%%%%%%%%%%%%%%%%%%%%%%%%%%%%%%%%%%%%%%%%%%%%%%%%%%%%%%%%%%%%%

%#begfold#{{{ Details for command \inptext
% \inptext[#1]{#2}
%   #1 = (escaped) pretext to typeset before #2
%   #2 = (escaped) text to typeset as an input variable.
% Typeset an input variable w/o hyperlink and index entry.
%#endfold#}}}
%
\NewDocumentCommand\inptext{O{}m}{%
  \inpstyle[#1]{#2}%
}

%%%%%%%%%%%%%%%%%%%%%%%%%%%%%%%%%%%%%%%%%%%%%%%%%%%%%%%%%%%%%%%%%%%%%%%%%%%

%#begfold#{{{ Details for command \inpindex
% \inpindex*{#2}[#3]
%   #1 = * (or not): prime index entry, else standard index entry
%   #2 = (escaped) text to \inpstyle for index entry
%   #3 = (escaped) text to append to index entry
%        Begin #3 with '!' to add a subentry.
% Called by other commands to create an input entry within the index.
% Examples:
%   \inpindex*{var\_name}[!channel]
%      => \index{Input!\texttt{var\_name}!channel|primeidx}
%   \inpindex{var\_name}[!channel]
%      => \index{Input!\texttt{var\_name}!channel}
%#endfold#}}}
%
\NewDocumentCommand\inpindex{smO{}}{%
  \IfBooleanTF#1
    {\index{Input!\inpstyle{#2}{#3}|primeidx}}% #1=star is present
    {\index{Input!\inpstyle{#2}{#3}}}%          #1=star is missing
}

%%%%%%%%%%%%%%%%%%%%%%%%%%%%%%%%%%%%%%%%%%%%%%%%%%%%%%%%%%%%%%%%%%%%%%%%%%%

%#begfold#{{{ Details for command \inp
% \inp*[#2]{#3}[#4][#5]    # Define/Reference an input variable.
%   #1 = * (or not): prime index entry, else standard index entry
%   #2 = hypertarget  (default=none)
%   #3 = (escaped) text to \inpstyle within paragraph
%   #4 = (escaped) text to \inpstyle for index entry (default=#3)
%   #5 = (escaped) text to append    to  index entry (default=none)
%        Begin #5 with '!' to add a subentry.
% - Used to define/refer-to an input variable.
% - Typeset the input variable.
% - Create/reference a hyperlink target.
% - Add the input variable to the index.
% Note:
%   Best to avoid special symbols in the hypertarget.
%   hypertarget with \_ gives undefined control sequence \textunderscore.
%   hypertarget with  _ is ok.
%#endfold#}}}
%
\NewDocumentCommand\inp{sO{}mO{=}O{}}{%
  \IfBooleanTF#1{% #1=star is present
    % Typeset the input variable, create optional hypertarget.
    \IfEq{#2}{}% Test if #2=hypertag is missing.
      {\inpstyle{#3}}%                             #2 is missing
      {\hypertarget{#2}{\inpstyle{#3}}%
       \phantomsection\label{#2}%
      }% #2 is present
      % pageref works with labels but not hypertargets.
      % phantomsection creates a local anchor, not the previous section header.
    % Typeset the variable in the index, adding optional sub-entry.
    \IfEq{#4}{=}% Test if #4=text to inpstyle for index is missing.
      {\inpindex*{#3}{#5}}% #4 is missing
      {\inpindex*{#4}{#5}}% #4 is present
  }{% #1=star is missing
    % Typeset the input variable, create optional hypertarget.
    \IfEq{#2}{}% Test if #2=hypertag is missing.
      {\inpstyle{#3}}%                 #2 is missing
      {\hyperlink{#2}{\inpstyle{#3}}}% #2 is present
    % Typeset the variable in the index, adding optional sub-entry.
    \IfEq{#4}{=}% Test if #4=text to inpstyle for index is missing.
      {\inpindex{#3}{#5}}% #4 is missing
      {\inpindex{#4}{#5}}% #4 is present
  }%
}

%%%%%%%%%%%%%%%%%%%%%%%%%%%%%%%%%%%%%%%%%%%%%%%%%%%%%%%%%%%%%%%%%%%%%%%%%%%

% Concat the left/right parts of a Fortran type variable: left%right.
%
\NewDocumentCommand{\concatstr}{mm}{%
  \IfEqCase{#2}{%
    {}{#1}%       If #2=''   return '#1'   w/o trailing '%'.
    {\%}{#1\%}%   If #2='\%' return '#1\%' w/  trailing '%'.
  }[#1\%#2]%      Else       return '#1\%#2'.
}

%##########################################################################
%##########################################################################
%##########################################################################
%##########################################################################
%##########################################################################
%  "inpBC" commands:
%   These are commands for creating index references that are for
%   exclusively for components of the BC_Input input derived type.
%   These are special versions of their respective "inp" commands but
%   help declutter the index by removing the BC_Input($i$)% prefix.
%   These commands also get the same special formatting/font styling
%   as the "inp" commands before placement in the index, 
%##########################################################################
%##########################################################################
%##########################################################################
%##########################################################################
%##########################################################################

%#begfold#{{{ Details for command \BCinpstyle
% \BCinpstyle[#1]{#2}
%   #1 = BC_Input index (default=i)
%   #2 = (escaped) text to typeset as a BC_Input component variable.
% Typeset a BC_Input variable w/o hyperlink and index entry.
%   Adds "BC\_Input(#1)\%" before the #2 variable name.
% Do not call directly.  Use \BCinptext instead.
%#endfold#}}}
% 
\NewDocumentCommand\BCinpstyle{O{i}m}{%
  {% Beg group to contain scope of commands.
%   \texttt{\MakeLowercase{\concatstr{BC\_Input($#1$)}{#2}}}%
%   \texttt{\MakeUppercase{\concatstr{BC\_Input($#1$)}{#2}}}%
                   \texttt{\concatstr{BC\_Input($#1$)}{#2}}%
%                         {\concatstr{BC\_Input($#1$)}{#2}}%
  }% End group to contain scope of commands.
}

%%%%%%%%%%%%%%%%%%%%%%%%%%%%%%%%%%%%%%%%%%%%%%%%%%%%%%%%%%%%%%%%%%%%%%%%%%%

%#begfold#{{{ Details for command \BCinptext
% \BCinptext[#1]{#2}
%   #1 = BC_Input index (default=i)
%   #2 = (escaped) text to typeset as an input variable.
% Called by other commands to apply the style for an input variable.
% Do not call directly.  Use \inptext or \BCinptext instead.
% The change case commands do not affect content in math mode.
%#endfold#}}}
\NewDocumentCommand\BCinptext{O{i}m}{%
  \BCinpstyle[#1]{#2}%
}

%%%%%%%%%%%%%%%%%%%%%%%%%%%%%%%%%%%%%%%%%%%%%%%%%%%%%%%%%%%%%%%%%%%%%%%%%%%

%#begfold#{{{ Details for command \inpBCindex (special version of \inpindex)
% SPECIAL VERSION OF \inpindex named \inpBCindex FOR COMPONENTS OF BC_INPUT_T.
% THIS IS USED TO DECLUTTER THE INDEX BY REMOVING THE BC_INPUT($i$)% PREFIX.
% \inpBCindex*[#2]{#3}[#4]
%   #1 = * (or not): prime index entry, else standard index entry
%   #2 = BC_Input index (default=i)
%   #3 = (escaped) text to \inpstyle for index entry
%   #4 = (escaped) text to append to index entry
%        Begin #3 with '!' to add a subentry.
% Used to create a BC_Input variable entry within the index.
% Examples:
%   \inpBCindex*{var\_name}[!channel]
%      => \index{Input!\texttt{BC\_Input($i$)%var\_name}!channel|primeidx}
%   \inpBCindex{var\_name}[!channel]
%      => \index{Input!\texttt{BC\_Input($i$)%var\_name}!channel}
%#endfold#}}}
%
\NewDocumentCommand\inpBCindex{sO{i}mO{}}{%
  % Always use $i$ instead of #2 when creating the index entry.
  \IfBooleanTF#1{% #1=star is present
    %#begfold#{{{ uncomment to dump debugging info
    %{ % Combine everything inside these curly brackets
    %  % onto a single line to dump debugging info.
    %  \color{Blue} 
    %  inpBCindex 2="#2",3="#3",4="#4":
    %  IfValue(3)=\IfValueTF{#3}{T}{F}
    %  "\inpstyle[\%]{#3}"
    %}
    %#endfold#}}}
     % The following nested logic is to prevent the index from adding
     % references to '%' nested under BC_INPUT when either \BCinpBC or
     % \inpBC is used without a BC_Input component variable in arguments
     % #4 or #5 to those commands.
     \IfEq{#4}{}% Test if #4 is missing
       % #4 is missing
       {\IfEq{#3}{}% Test if #3 is empty
         % #3 == '', return without %
         {\index{Input!\inpstyle{BC\_Input($i$)}|primeidx}}%
         % #3 \= '', return with \inpstyle{%#3}
         {\index{Input!\inpstyle{BC\_Input($i$)}!{\inpstyle[\%]{#3}}|primeidx}}%
       }%
       % #4 is present
       {\IfEq{#3}{}% Test if #3 is empty
         % #3 == '', return with \inpstyle{%}#4
         {\index{Input!\inpstyle{BC\_Input($i$)}!{\inpstyle[\%]{}}{#4}|primeidx}}%
         % #3 \= '', return with \inpstyle{%#3}#4
         {\index{Input!\inpstyle{BC\_Input($i$)}!{\inpstyle[\%]{#3}}{#4}|primeidx}}%
       }%
    % THE FOLLOWING LINE IS THE OLD CALL TO THE INDEX FUNCTION
    % THAT RESULTED IN SEPARATE '%' REFERENCES UNDER BC_INPUT
  %% \index{Input!\inpstyle{BC\_Input($i$)}!{\inpstyle[\%]{#3}}{#4}|primeidx}%
    } {% #1=star is missing
     % The following nested logic is to prevent the index from adding
     % references to '%' nested under BC_INPUT when either \BCinpBC or
     % \inpBC is used without a BC_Input component variable in arguments
     % #4 or #5 to those commands.
     \IfEq{#4}{}% Test if #4 is missing
       % #4 is missing
       {\IfEq{#3}{}% Test if #3 is empty
         % #3 == '', return without %
         {\index{Input!\inpstyle{BC\_Input($i$)}}}%
         % #3 \= '', return with \inpstyle{%#3}
         {\index{Input!\inpstyle{BC\_Input($i$)}!{\inpstyle[\%]{#3}}}}%
       }%
       % #4 is present
       {\IfEq{#3}{}% Test if #3 is empty
         % #3 == '', return with \inpstyle{%}#4
         {\index{Input!\inpstyle{BC\_Input($i$)}!{\inpstyle[\%]{}}{#4}}}%
         % #3 \= '', return with \inpstyle{%#3}#4
         {\index{Input!\inpstyle{BC\_Input($i$)}!{\inpstyle[\%]{#3}}{#4}}}%
       }%
    % THE FOLLOWING LINE IS THE OLD CALL TO THE INDEX FUNCTION
    % THAT RESULTED IN SEPARATE '%' REFERENCES UNDER BC_INPUT
  %% \index{Input!\inpstyle{BC\_Input($i$)}!{\inpstyle[\%]{#3}}{#4}}%
    }%
}

%%%%%%%%%%%%%%%%%%%%%%%%%%%%%%%%%%%%%%%%%%%%%%%%%%%%%%%%%%%%%%%%%%%%%%%%%%%

%#begfold#{{{ Details for command \inpBC (special version of \inp)
% SPECIAL VERSION OF \inp named \inpBC FOR COMPONENTS OF BC_INPUT_T.
% ADD NO PREFIX TO DOCUMENT ENTRIES.
% ADD PREFIX "BC_INPUT(#3)%" TO INDEX ENTRIES.
% \BCinp*[#2][#3]{#4}[#5][#6]    # Define/Reference a BC_Input sub-variable.
%   #1 = * (or not): prime index entry, else standard index entry
%   #2 = hypertarget  (default=none)
%   #3 = BC_Input index (default=i)
%   #4 = (escaped) text to \inpstyle within paragraph
%   #5 = (escaped) text to \inpstyle for index entry (default=#4)
%        Will be prefixed by "BC\_Input(#3)%".
%   #6 = (escaped) text to append    to  index entry (default=none)
%        Begin #6 with '!' to add a subentry.
% - Used to define/refer-to an input variable within the document.
% - Typeset the input variable.
% - Create/reference the hyperlink target.
% - Add the input variable to the index.
% Note:
%   Best to avoid special symbols in the hypertarget.
%   hypertarget with \_ gives undefined control sequence \textunderscore.
%   hypertarget with  _ is ok.
%#endfold#}}}
%
\NewDocumentCommand\inpBC{sO{}O{i}mO{=}O{}}{%
  \IfBooleanTF#1{% #1=star is present
    % Typeset the input variable in the document, create optional hypertarget.
    \IfEq{#2}{}% Test if #2=hypertag is missing.
      {\inpstyle{#4}}%                   #2 is missing
      {\hypertarget{#2}{\inpstyle{#4}}}% #2 is present
    % Typeset the input variable in the index, adding optional sub-entry.
    \IfEq{#5}{=}% Test if #5=text to style for index is missing.
      {\inpBCindex*[#3]{#4}{#6}}% #5 is missing
      {\inpBCindex*[#3]{#5}{#6}}% #5 is present
  }{% #1=star is missing
    % Typeset the input variable in the document, create optional hyperlink.
    \IfEq{#2}{}% Test if #2=hypertag is missing.
      {\inpstyle{#4}}%                 #2 is missing
      {\hyperlink{#2}{\inpstyle{#4}}}% #2 is present
    % Typeset the input variable in the index, adding optional sub-entry.
    \IfEq{#5}{=}% Test if #5=text to inpBCstyle for index is missing.
      {\inpBCindex[#3]{#4}{#6}}% #5 is missing
      {\inpBCindex[#3]{#5}{#6}}% #5 is present
  }%
}

%%%%%%%%%%%%%%%%%%%%%%%%%%%%%%%%%%%%%%%%%%%%%%%%%%%%%%%%%%%%%%%%%%%%%%%%%%%

%#begfold#{{{ Details for command \BCinpBC 
% SPECIAL VERSION OF \inp named \BCinpBC FOR COMPONENTS OF BC_INPUT_T.
% ADD PREFIX "BC_INPUT(#3)%" TO DOCUMENT ENTRIES.
% ADD PREFIX "BC_INPUT(#3)%" TO INDEX ENTRIES.
% \BCinp*[#2][#3]{#4}[#5][#6]    # Define/Reference a BC_Input sub-variable.
%   #1 = * (or not): prime index entry, else standard index entry
%   #2 = hypertarget  (default=none)
%   #3 = BC_Input index (default=i)
%   #4 = (escaped) text to \inpstyle within paragraph
%   #5 = (escaped) text to \inpstyle for index entry (default=#4)
%   #6 = (escaped) text to append    to  index entry (default=none)
%        Begin #6 with '!' to add a subentry.
% - Used to define/refer-to an input variable within the document.
% - Typeset the input variable.
% - Create/reference the hyperlink target.
% - Add the input variable to the index.
% Note:
%   Best to avoid special symbols in the hypertarget.
%   hypertarget with \_ gives undefined control sequence \textunderscore.
%   hypertarget with  _ is ok.
%#endfold#}}}
%
\NewDocumentCommand\BCinpBC{sO{}O{i}mO{=}O{}}{%
  \IfBooleanTF#1{% #1=star is present
    % Typeset the input variable in the document, create optional hypertarget.
    \IfEq{#2}{}% Test if #2=hypertag is missing.
      {\BCinpstyle[#3]{#4}}%                   #2 is missing
      {\hypertarget{#2}{\BCinpstyle[#3]{#4}}}% #2 is present
    % Typeset the input variable in the index, adding optional sub-entry.
    \IfEq{#5}{=}% Test if #5=text to style for index is missing.
      {\inpBCindex*[#3]{#4}{#6}}% #5 is missing
      {\inpBCindex*[#3]{#5}{#6}}% #5 is present
  }{% #1=star is missing
    % Typeset the input variable in the document, create optional hyperlink.
    \IfEq{#2}{}% Test if #2=hypertag is missing.
      {\BCinpstyle[#3]{#4}}%                   #2 is missing
      {\hyperlink{#2}{\BCinpstyle[#3]{#4}}}% #2 is present
    % Typeset the input variable in the index, adding optional sub-entry.
    \IfEq{#5}{=}% Test if #5=text to inpBCstyle for index is missing.
      {\inpBCindex[#3]{#4}{#6}}% #5 is missing
      {\inpBCindex[#3]{#5}{#6}}% #5 is present
  }%
}

%%%%%%%%%%%%%%%%%%%%%%%%%%%%%%%%%%%%%%%%%%%%%%%%%%%%%%%%%%%%%%%%%%%%%%%%%%%

\NewDocumentEnvironment{nmldescbody}{}%
  % Code to execute before begin:
  {\vspace{-0.6\baselineskip}\begin{adjustwidth}{\parindent}{}}
  % Code to execute before end:
  {\end{adjustwidth}\vspace{0.2\baselineskip}}

%%%%%%%%%%%%%%%%%%%%%%%%%%%%%%%%%%%%%%%%%%%%%%%%%%%%%%%%%%%%%%%%%%%%%%%%%%%

% SI Units. Use these where possible for consistency.
  \newcommand{\unitPres}{\unit{\Pa}}
  \newcommand{\unitTemp}{\unit{\K}}
  \newcommand{\unitDens}{\unit{\kg\per\m\cubed}}
  \newcommand{\unitVel }{\unit{\m\per\s}}

%%%%%%%%%%%%%%%%%%%%%%%%%%%%%%%%%%%%%%%%%%%%%%%%%%%%%%%%%%%%%%%%%%%%%%%%%%%


% ======================
% Math / Equation Macros
% ======================

% Declare a math small cap font.
  \DeclareMathAlphabet{\mathsc}{OT1}{cmr}{m}{sc}

% r-variable, coord in the standard element [-1:1].
  \newcommand{\rv}{\ensuremath{r}} % r, coord in elem [-1:1]

% Parenthesis-Variable-Parenthesis: Put parentheses around a variable.
%   \pvp    = (r)
%   \pvp[x] = (x)
  \newcommand{\pvp}[1][\rv]{\ensuremath{\!\left(#1\right)}}

% Symbol to represent the standard element domain [-1:1].
  \newcommand{\stdelem}{\ensuremath{\mathbf{I}}} % std elem domain [-1:1]

% Nondimensional time is (time * Time_Scaling_Factor), denoted by: t^*.
  \newcommand{\tstar}{\ensuremath{t^{*}}}

  \newcommand{\solvar}{u}                    % solution variable
  \newcommand{\flxvar}{f}                    % flux     variable
  \newcommand{\solfun}{\ensuremath{\solvar}} % solution function
  \newcommand{\flxfun}{\ensuremath{\flxvar}} % flux     function

% Polynomial space.
% Seth's commands for polynomial order and cutoff order.
% ToDo: Need a comment here on how to use these!
% ToDo: Combine with \PO command below.
  \newcommand{\p}[2][]{\ensuremath{\mathcal{P}_{\textsl{#2}}#1}}
    % \p[trailing]{subscript}
    %   Note "trailing" text is typeset AFTER the subscript.
    %   [trailing]  is often [3] => P3, [+1] => P+1
    %   {subscript} is often P=projection, S=solution, Q=quadrature
    % \p{3}      =>  \mathcal{P}_{\textsl{3}    (ie, P_3                   )
    % \p[+1]{}   =>  \mathcal{P}_{\textsl{} +1  (ie, P+1                   )
    % \p[+1]{S}  =>  \mathcal{P}_{\textsl{S}+1  (ie, P_S+1 where S=solution)
  \newcommand{\pc}[2][]{\ensuremath{\mathcal{P}c#1_{\textsl{#2}}}}
    % \pc[trailing]{subscript}
    %   Note "trailing" text is typeset BEFORE the subscript.

% Polynomial order and modes.
% Originally hard-coded in tex. Need to go back and replace with macros.
% used: P   =      polynomial order
% used: N_p = P+1  polynomial modes
% \newcommand{\PO}[1][]{\ensuremath{P#1}} % polynomial order, P
  \newcommand{\PO}[1][]{\ensuremath{{\cal P}#1}} % polynomial order, P
% \newcommand{\PO}[1][]{\ensuremath{\mathbb{P}#1}} % polynomial order, P
    % #1 = specific polynomial order (ie, 0-6), not an offset (ie, +1)
    % \PO     =>  P
    % \PO[3]  =>  P3
    % \PO+1   =>  P+1
% \newcommand{\PM}{\ensuremath{N_p}} % polynomial modes, N_p = P+1
  \newcommand{\PM}{\ensuremath{N_{{\cal P}}}} % polynomial modes, N_p = P+1
% \newcommand{\PM}{\ensuremath{N_\mathbb{P}}} % polynomial modes, N_p = P+1

% nodal and modal solution variables
%   Optional arg#1 = subscript
  \newcommand{\uNodal}[1][]{\ensuremath{\solfun_{#1}}}
  \newcommand{\uModal}[1][]{\ensuremath{\hat{\solfun}_{#1}}}

  \newcommand{\LB}{\ensuremath{\mathsc{lb}}} % left  boundary (of cell)
  \newcommand{\RB}{\ensuremath{\mathsc{rb}}} % right boundary (of cell)

% Upwind flux functions
  \newcommand{\fupw}[1][]{\ensuremath{\flxfun^{\ast}_{#1}}} % upwind flux
    % Optional arg#1 = subscript
  \newcommand{\fupwL}{\ensuremath{\fupw[\LB]}}  % upwind flux left
  \newcommand{\fupwR}{\ensuremath{\fupw[\RB]}}  % upwind flux right
  \newcommand{\fC}{\ensuremath{F}}              % continuous flux function
  \newcommand{\fD}[1][]{\ensuremath{f^D_{#1}\mspace{-2.0mu}}} % discont flux fn
    % Optional arg#1 = subscript

% left/right solution variables for flux calculation across face
  \newcommand{\uLeft}{\ensuremath{\solfun_{\mathsc{l}}}}
  \newcommand{\uRight}{\ensuremath{\solfun_{\mathsc{r}}}}

% left/right flux jumps. Flux jumps on the left/right edge of cell.
  \newcommand{\fluxjump}[1][]{\ensuremath{\Delta f_{#1}}}    % flux jump
  \newcommand{\fluxjumpL}{\ensuremath{\fluxjump[\LB]}} % l-bnd flux jump
  \newcommand{\fluxjumpR}{\ensuremath{\fluxjump[\RB]}} % r-bnd flux jump

% unit correction function for fluxes
  \newcommand{\g}[1][]{\ensuremath{\text{\textsl{g}}_{#1}}} % unit corr fun
  \newcommand{\gL}{\ensuremath{\g[\LB]}}         % l-boundary unit corr fun
  \newcommand{\gR}{\ensuremath{\g[\RB]}}         % r-boundary unit corr fun

% scaled correction function for fluxes
  \newcommand{\G}[1][]{\ensuremath{G_{#1}}} % unit corr fun
  \newcommand{\GL}{\ensuremath{\G[\LB]}}         % l-boundary unit corr fun
  \newcommand{\GR}{\ensuremath{\G[\RB]}}         % r-boundary unit corr fun

% Legendre polynomial
  \newcommand{\Leg}[1]{\ensuremath{P_{#1}}}              % Legendre
  \newcommand{\dLeg}[1]{\ensuremath{P^{\,\prime}_{#1}}}  % Legendre deriv

% Radau polynomial
  \newcommand{\RadL}[1]{\ensuremath{R_{\mathsc{L},#1}}}  % Radau Left
  \newcommand{\RadR}[1]{\ensuremath{R_{\mathsc{R},#1}}}  % Radau Right
  \newcommand{\dRadL}[1]{\ensuremath{R^{\,\prime}_{\mathsc{L},#1}}} % RadL deriv
  \newcommand{\dRadR}[1]{\ensuremath{R^{\,\prime}_{\mathsc{R},#1}}} % RadR deriv


% Vandermonde matrix
% \newcommand{\vand}{\ensuremath{\bm{\cal V}}} % fails
  \newcommand{\vand}{\ensuremath{   {\cal V}}}

% Levi-Civita permutation tensor
  \newcommand{\perm}[1]{\ensuremath{\varepsilon_{#1}}}

% unit vector
  \newcommand{\bmunitvec}[1]{\ensuremath{\hat{\bm{e}}_{{#1}}}} % unit vector
  \newcommand{\unitvec}[1]{\ensuremath{\hat{e}_{{#1}}}} % unit vector

%
% Commands to typeset vectors, tensors, and matrix variables
% in symbol (display) and index forms.
%
%
% Vector symbol. Optional arg #1 is accent over symbol.
% Usage: \vectorsymbol[\hat]{x}
  \newcommand{\vectorsymbolvec  }[2][]{\vec{#1{#2}}}
  \newcommand{\vectorsymbololine}[2][]{\overline{#1{#2}}}
  \newcommand{\vectorsymboluline}[2][]{\underline{#1{#2}}}
  \newcommand{\vectorsymbolbold }[2][]{#1{\bm{{#2}}}}
  \newcommand{\vectorsymbol     }[2][]{\vectorsymbolvec[#1]{#2}}
%
% Tensor symbol. Optional arg #1 is accent over symbol.
% Usage: \tensorsymbol[\hat]{\tau}
  \newcommand{\tensorsymbolvec  }[2][]{\vec{\vec{#1{#2}}}}
  \newcommand{\tensorsymbololine}[2][]{\overline{\overline{#1{#2}}}}
  \newcommand{\tensorsymboluline}[2][]{\underline{\underline{#1{#2}}}}
  \newcommand{\tensorsymbolbold }[2][]{#1{\bm{{#2}}}}
  \newcommand{\tensorsymbol     }[2][]{\tensorsymbolbold[#1]{{#2}}}
%
% Matrix symbol. Optional arg #1 is accent over symbol.
% Usage: \matrixsymbol[\hat]{Q}
  \newcommand{\matrixsymbololine}[2][]{\overline{\overline{#1{#2}}}}
  \newcommand{\matrixsymboluline}[2][]{\underline{\underline{#1{#2}}}}
  \newcommand{\matrixsymbolbold }[2][]{#1{\bm{{#2}}}}
  \newcommand{\matrixsymbolsbrac}[2][]{\left[#1{\bm{#2}}\right]}
  \newcommand{\matrixsymbol     }[2][]{\matrixsymbolbold[#1]{#2}}
%
%
% Vector index. Optional arg #1 is accent over symbol.
% Usage: \vectorindex[\hat]{x}{i}
  \newcommand{\vectorindex}[3][]{#1{{#2}}_{#3}}
%
% Tensor index. Optional arg #1 is accent over symbol.
% Usage: \tensorindex[\hat]{\tau}{ij}
  \newcommand{\tensorindex}[3][]{#1{{#2}}_{#3}}
%
% matrix index. Optional arg #1 is accent over symbol.
% Usage: \matrixindex[\hat]{Q}{ij}
  \newcommand{\matrixindex}[3][]{#1{{#2}}_{#3}}


%
% Rotating frame notation.
% Mixed frame means a fixed frame quantity written in terms of the
% rotating frame unit vectors.
%
  \newcommand{\fix}[1]{{#1}}               % plain symbol   in fixed    frame
  \newcommand{\rot}[1]{{\hat{#1}}}         % plain symbol   in rotating frame
  \newcommand{\mix}[1]{{\acute{#1}}}       % plain symbol   in mixed    frame
% Use this form to note objective scalars, where fix and rot have same value.
% Can then easily change whether those are typeset as \rot or \fix.
  \newcommand{\obj}[1]{{\fix{#1}}}         % plain scalar   in rotating frame
% \newcommand{\obj}[1]{{\hat{#1}}}         % plain scalar   in rotating frame


%
% Specialized commands for fix, rot, mix frame
% to typeset vectors, tensors, and matrix variables
% in symbol (display) and index forms.
%
%
% Vector symbols for fix, rot, mix frame.
% Usage: \fvs{x}{i}
  \newcommand{\fvs}[2][]{\vectorsymbolvec[#1]{{\fix{#2}}}}
  \newcommand{\rvs}[2][]{\vectorsymbolvec[#1]{{\rot{#2}}}}
  \newcommand{\mvs}[2][]{\vectorsymbolvec[#1]{{\mix{#2}}}}
%
% Tensor symbols for fix, rot, mix frame.
% Usage: \fts{\tau}{ij}
  \newcommand{\fts}[2][]{\tensorsymbolbold[#1]{{\fix{#2}}}}
  \newcommand{\rts}[2][]{\tensorsymbolbold[#1]{{\rot{#2}}}}
  \newcommand{\mts}[2][]{\tensorsymbolbold[#1]{{\mix{#2}}}}
%
% Matrix symbols for fix, rot, mix frame.
% Usage: \fms{Q}{ij}
  \newcommand{\fms}[2][]{\matrixsymbolbold[#1]{{\fix{#2}}}}
  \newcommand{\rms}[2][]{\matrixsymbolbold[#1]{{\rot{#2}}}}
  \newcommand{\mms}[2][]{\matrixsymbolbold[#1]{{\mix{#2}}}}
%
%
% Vector index for fix, rot, mix frame.
% Usage: \fvi[\dot]{x}{i}
  \newcommand{\fvi}[3][]{\vectorindex[#1]{{\fix{#2}}}{#3}}
  \newcommand{\rvi}[3][]{\vectorindex[#1]{{\rot{#2}}}{#3}}
  \newcommand{\mvi}[3][]{\vectorindex[#1]{{\mix{#2}}}{#3}}
%
% Tensor index for fix, rot, mix frame.
% Usage: \fti[\dot]{\tau}{ij}
  \newcommand{\fti}[3][]{\tensorindex[#1]{{\fix{#2}}}{#3}}
  \newcommand{\rti}[3][]{\tensorindex[#1]{{\rot{#2}}}{#3}}
  \newcommand{\mti}[3][]{\tensorindex[#1]{{\mix{#2}}}{#3}}
%
% Matrix index for fix, rot, mix frame.
% Usage: \fmi[\dot]{Q}{ij}
  \newcommand{\fmi}[3][]{\matrixindex[#1]{{\fix{#2}}}{#3}}
  \newcommand{\rmi}[3][]{\matrixindex[#1]{{\rot{#2}}}{#3}}
  \newcommand{\mmi}[3][]{\matrixindex[#1]{{\mix{#2}}}{#3}}


%
% More commands for rotating reference frame.
%
% Rotation axis components in fixed frame.
  \newcommand{\fvia}[1]{\fvi{a}{\fix{#1}}}
% Rotation axis components in rotating frame.
  \newcommand{\mvia}[1]{\mvi{a}{\rot{#1}}}
% Rotating frame origin translational offset vector variable.
  \newcommand{\FrameTransVecVar}{o}
% Rotating frame spin vector variable.
  \newcommand{\FrameSpinVecVar}{\omega}
% Rotating frame spin tensor variable.
  \newcommand{\FrameSpinTenVar}{\Omega}
% Rotation matrix to change from fixed to rotated frame.
  \newcommand{\FrameRotMatVar}{Q}
% Rotation angle and rotation rate.
  \newcommand{\FrameRotAngle}{\ensuremath{\theta}}
  \newcommand{\FrameRotRate}{\ensuremath{\dot{\theta}}}
% Frame velocity.
  \newcommand{\FrameVelVar}{w}


%%The following are no longer used and can be removed.
%%Retained here until we delete the older backup .tex files.
%
% \newcommand{\bmfix}[1]{{\bm #1}}         % bold  symbol   in fixed    frame
% \newcommand{\bmrot}[1]{{\hat{\bm{#1}}}}  % bold  symbol   in rotating frame
% \newcommand{\bmmix}[1]{{\check{\bm{#1}}}}% bold  symbol   in mixed    frame
%
% \newcommand{\vecfix}[1]{\vec{#1}}        % vect  symbol   in fixed    frame
% \newcommand{\vecrot}[1]{\vec{\hat{#1}}}  % vect  symbol   in rotating frame
% \newcommand{\vecmix}[1]{\vec{\check{#1}}}% vect  symbol   in mixed    frame
%
% \newcommand{      \FrameVel}{w}             % velocity var   of rotating frame
% \newcommand{   \rotFrameVel}{\hat{w}}       % plain velocity of rotating frame
% \newcommand{ \bmrotFrameVel}{\hat{\bm w}}   % bold  velocity of rotating frame
% \newcommand{\vecrotFrameVel}{\vec{\hat{w}}} % vect  velocity of rotating frame


% Markers.
  \newcommand{\marker}{\ensuremath{{\cal M}}}
  \newcommand{\scaledmarker}{\ensuremath{{\cal S}}}

% Lengths used in description environment.
  \newlength{\itemIndentLength}
  \setlength{\itemIndentLength}{1.5ex}
  \newlength{\descLabelLength}
  \setlength{\descLabelLength}{\widthof{$\bullet$\hspace{\itemIndentLength}}}

% FIXME command moved to latex/gfr.sty

% QandA environment: 1 optional argument for color of questions (default black)
  \NewDocumentEnvironment{QandA}{O{black}}
    {\begin{description}[align=left,
                         style=nextline,
                         font=\color{#1}$\bullet$\hspace{\itemIndentLength},
                         leftmargin=\descLabelLength]}
    {\end{description}}

% InputList environment
\NewDocumentEnvironment{InputList}{O{}}
% Avoid font color and style, since content is often a user input variable
%   with its own font style and hyperlink color.
% {\begin{itemize}[align=left,style=nextline,
%                  font=\color{blue}\slshape,
%                  leftmargin=\descLabelLength,labelsep=40em]}
  {\begin{description}[align=left,
                       style=nextline,
                       leftmargin=\descLabelLength
                       ,#1]}
%                      \IfValueTF{#1}{#1}{}]}
  {\end{description}}

% Custom itemize environment used for formatting a group
% of input variables and respective values as an example
\NewDocumentEnvironment{InputExample}{O{}}
  {\begin{itemize}[left=3.0em,
                   align=left,
                   style=nextline,
                   topsep=3pt,
                   noitemsep,
                   #1]}
  {\end{itemize}}

% Renaming paragraph as subsubsubsection
% and aliased to ThreeSubsection
\newlength{\tmpParagraphLength}
\NewDocumentCommand{\subsubsubsection}{moo}{
  \IfValueTF{#3}{
    \setlength{\tmpParagraphLength}{#3}
  } {
    \setlength{\tmpParagraphLength}{0pt}
  }
  \IfValueTF{#2}{
   %\paragraph{#1}\label{#2}{}\ {}\\[\tmpParagraphLength]
    \paragraph{#1}\label{#2}
  } {
   %\paragraph{#1}{}\ {}\\[\tmpParagraphLength]
    \paragraph{#1}
  }
}
\NewCommandCopy{\ThreeSubsection}{\subsubsubsection}

% Renaming subparagraph as subsubsubsubsection
% and aliased to FourSubsection
\NewDocumentCommand{\subsubsubsubsection}{moo}{
  \IfValueTF{#3}{
    \setlength{\tmpParagraphLength}{#3}
  } {
    \setlength{\tmpParagraphLength}{0pt}
  }
  \IfValueTF{#2}{
   %\subparagraph{#1}\label{#2}{}\ {}\\[\tmpParagraphLength]
    \subparagraph{#1}\label{#2}
  } {
   %\subparagraph{#1}{}\ {}\\[\tmpParagraphLength]
    \subparagraph{#1}
  }
}
\NewCommandCopy{\FourSubsection}{\subsubsubsubsection}

% Use this wrapper for images, in case we decide to add a box later.
\NewDocumentCommand{\insertimage}{om}{
  \IfValueTF{#1}{
    \includegraphics[#1]{#2}
  } {
    \includegraphics{#2}
  }
}

% Use this wrapper for subfigure environment, 
% also for in case we decide to add a box later.
\NewDocumentCommand{\subfig}{mm}{
  \begin{subfigure}[t]{#1}\centering#2\end{subfigure}
}

% Various lengths used throughout for sizing figures and minipages
\newlength{\imagewidth}
\newlength{\imageheight}
\newlength{\mpwidth}
\newlength{\mpheight}
%\newlength{\capParSep} \setlength{\capParSep}{4pt}
\newlength{\capParSep} \setlength{\capParSep}{1.0ex}

%\renewcommand\LTcaptype{}
% Older versions of longtable package give error from:
%  \renewcommand\LTcaptype{}
%  LaTeX Error: No counter '' defined.
% so give it a dummy counter that is not used.
\newcounter{dummycounterlongtable}
\renewcommand\LTcaptype{dummycounterlongtable}
